%  Wstęp - wprowadzenie w problem/zagadnienia, osadzenie problemu w dziedzinie, cel pracy,
% zakres pracy, zwięzła charakterystyka rozdziałów, jednoznaczne określenie wkładu autora.
\chapter{Wstęp}
\label{c:wstep}

\section{Motywacja podjęcia tematu}

Utrudnienia codziennego funkcjonowania osób niewidomych i słabowidzących zauważyć można każdego dnia. Wiele czynności, które dla osób widzących są proste i intuicyjne, dla osób z dysfunkcją wzroku wymagają zastosowania specjalnych metod i narzędzi. Korzystanie z transportu publicznego, płatności gotówkowe czy też zwykłe rozpoznawanie przedmiotów codziennego użytku stanowią dla nich poważne wyzwanie. Realia szkolne i zawodowe również nie są wolne od trudności. Nauka abstrakcyjnych pojęć matematycznych, szczególnie związanych z geometrią przestrzenną, wymaga od uczniów stosowania specjalnych metod nauczania, które ułatwiają im zrozumienie i przyswojenie wiedzy. 

Przy nauce stereometrii, szczególnie użyteczne okazują się fizyczne modele brył, które dają możliwość poznania ich kształtu i właściwości poprzez dotyk. Jednakże same modele, mimo że umożliwiają dotykowe poznanie bryły, nie przekazują wszystkich informacji potrzebnych w procesie nauki. Często, aby uzyskać pełniejszy obraz, konieczne jest korzystanie z dodatkowych treści edukacyjnych, jak np. wydruki brajlowskie. Ich użycie może prowadzić jednak do przerwania procesu eksploracji i wydłużenia czasu pozyskiwania informacji, ponieważ wymaga od ucznia oderwania rąk od modelu i sięgnięcia po dodatkowe materiały.

Istniejące rozwiązania wspomagające osoby z dysfunkcją wzroku koncentrują się głównie na wspomaganiu użytkownika w codziennych czynnościach i przemieszczaniu się w przestrzeni. Istnieje jednak niewiele rozwiązań wspomagających naukę geometrii przestrzennej. Ten fakt stanowił główną motywację podjęcia tematu pracy -- stworzenia rozwiązania, które wesprze osoby niewidome i słabowidzące w nauce geometrii przestrzennej i ułatwi im dostęp do informacji o poznawanych obiektach. 

\section{Cel i zakres pracy}

Celem pracy jest opracowanie systemu wspomagającego osoby niewidome i słabowidzące w nauce geometrii przestrzennej za pomocą fizycznych modeli 3D, wykorzystującego metody detekcji obiektów na obrazie. Część badawcza pracy obejmuje przygotowanie i eksperymentalną ocenę użycia modeli uczenia głębokiego do wykrywania obiektów, skupiając się na ich skuteczności detekcji i czasie inferencji. W ramach analizy należy również ocenić wpływ wybranych parametrów treningu na wyniki detekcji. 

Zakres pracy obejmuje przeprowadzenie analizy literatury, przygotowanie działającego prototypu systemu, własnego zbioru danych do trenowania modeli detekcji, wykonanie eksperymentów w celu oceny skuteczności wybranych modeli oraz opracowanie wniosków i rekomendacji dotyczących dalszego rozwoju systemu. Praca nie skupia się na opracowaniu nowych metod detekcji obiektów, a jedynie na dostosowaniu istniejących rozwiązań do wymagań projektowanego systemu. Zakres eksperymentalny obejmuje ocenę wpływu rozmiaru modelu, rozdzielczości obrazu wejściowego oraz reprezentacji obrazu na skuteczność detekcji. 

W zakresie pracy nie przewidziano badań z udziałem osób niewidomych i słabowidzących, objęcia całego zakresu nauczania stereometrii ani oceny wszystkich dostępnych metod detekcji obiektów. 

\section{Wkład autora}

Wkład własny obejmuje część projektowo-implementacyjną oraz badawczą pracy. W ramach pierwszej z nich opracowano koncepcję i architekturę systemu oraz zaimplementowano aplikację mobilną i część serwerową. Przygotowano dla nich mechanizmy komunikacji i przesyłania danych, detekcji obiektów, ich analizy oraz generowania komunikatów głosowych. Opracowano również zestaw narzędzi pomocniczych do przygotowania danych, trenowania modeli i prowadzenia eksperymentów. 

Część badawcza objęła przygotowanie własnego zbioru danych, opracowanie metodyki eksperymentów oraz trening i ocenę siedmiu modeli detekcji obiektów. Zbadano wpływ skali modelu, rozdzielczości obrazu wejściowego oraz jego reprezentacji na skuteczność detekcji. Następnie wyniki poddano analizie i wskazano konfigurację modelu, która najlepiej odpowiada przyjętym wymaganiom. 

\section{Struktura pracy}

Pełny tekst pracy został podzielony na siedem rozdziałów. Każdy z nich obejmuje odrębny obszar tematyczny, a ich kolejność odpowiada przyjętej logice realizacji projektu.

W rozdziale \ref{c:analiza} przedstawiono analizę literatury związanej z tematem pracy. Omówiono w nim sposoby poznawania obiektów przez osoby niewidome i słabowidzące, rozwój technologii asystujących oraz problem rozpoznawania obiektów przy interakcji dotykowej. Następnie przedstawiono przegląd metod detekcji obiektów, stosowane metryki oceny, rodzinę modeli YOLO. Rozdział zakończony jest podsumowaniem zawierającym wnioski z analizy literatury wykorzystywane w dalszych etapach pracy.

Rozdział \ref{c:koncepcja} przedstawia koncepcję działania i architekturę systemu. Omówiono w nim wymagania, sposób przepływu danych, funkcjonalności oraz ograniczenia systemu. Szczegóły implementacyjne zostały przedstawione w rozdziale \ref{c:implementacja}, który opisuje proces implementacji aplikacji mobilnej i części serwerowej, sposób ich komunikacji oraz przygotowane narzędzia pomocnicze. 

Część badawczą pracy rozpoczyna rozdział \ref{c:zbior}, w którym opisano sposób przygotowania własnego zbioru danych oraz jego charakterystykę. Zawarto w nim również metodykę eksperymentów. W rozdziale \ref{c:badania} zostały przedstawione wyniki treningów i testów modeli. Porównano w nim wpływ czynników na skuteczność detekcji oraz przeanalizowano błędy popełniane przez modele. 

Podsumowanie pracy i wnioski z przeprowadzonych badań zawarto w rozdziale \ref{c:podsumowanie}. Wskazano w nim najważniejsze wnioski, możliwe kierunki rozwoju oraz ocenę stopnia realizacji celu pracy.
